\chapter{کلیات تحقیق}
\label{ch:introduction}

%==============================================================================
\section{بستر بالینی}
%==============================================================================

\subsection{محیط بخش مراقبت ویژه و تصمیم‌گیری}

بخش مراقبت‌های ویژه (\lr{ICU}) محیطی پرخطر و زمان‌حساس است؛ جایی که تصمیم‌های بالینی باید با اتکا به داده‌های ناقص اما پیوسته به‌روز گرفته شوند. بیماران این بخش معمولاً وضعیت ناپایدارتری دارند و تحت پایش مداوم علائم حیاتی، آزمایش‌های مکرر و مراقبت تیم‌های چندرشته‌ای قرار می‌گیرند. تصمیم‌هایی که در \pnum{24} تا \pnum{48} ساعت نخست بستری گرفته می‌شود، اغلب مسیر کل بستری را تعیین می‌کند؛ از تحرک‌دهی زودهنگام و جداسازی از ونتیلاسیون مکانیکی تا برنامه‌ریزی ترخیص، همگی به برآورد دقیق پایداری بالینی و ظرفیت بهبود بیمار وابسته‌اند.

\subsection{گردش کار بالینی و نقاط تصمیم}

گردش کار معمول \lr{ICU} را می‌توان در چهار گام خلاصه کرد: (\pnum{1}) پذیرش و ارزیابی اولیه؛ (\pnum{2}) تثبیت وضعیت در \pnum{24} تا \pnum{48} ساعت نخست؛ (\pnum{3}) پایش و درمان مداوم؛ و (\pnum{4}) برنامه‌ریزی ترخیص. ارزیابی امکان‌پذیری درمان معمولاً در مرحلۀ تثبیت انجام می‌شود، زمانی که نخستین روندهای آزمایشگاهی و علائم حیاتی در دسترس قرار گرفته‌اند. بالین‌گران در این مرحله چند جریان داده را کنار هم می‌گذارند: پانل‌های آزمایشگاهی مانند کراتینین، لاکتات و گلوکز؛ علائم حیاتی مانند ضربان قلب، فشار خون و \lr{SpO$_2$}؛ و یادداشت‌های بالینی. مدلی که بتواند این داده‌ها را به‌صورت منسجم ترکیب کند، می‌تواند بیمارانی را که شانس ترخیص ایمن زودهنگام دارند زودتر شناسایی کند و برنامه‌ریزی ترخیص را دقیق‌تر سازد.

\subsection{برنامه‌ریزی ترخیص و امکان‌پذیری درمان}

برنامه‌ریزی ترخیص از آغاز بستری در \lr{ICU} آغاز می‌شود. بالین‌گران باید به‌طور مداوم ارزیابی کنند که آیا بیمار می‌تواند به‌صورت ایمن به منزل، توان‌بخشی یا مرکز دیگری ترخیص شود. این ارزیابی (\emph{امکان‌پذیری درمان})، روندهای آزمایشگاهی (مثلاً بهبود کراتینین در نارسایی کلیوی)، پایداری علائم حیاتی (مثلاً جداسازی از وازوپرسورها)، وضعیت عملکردی و عوامل اجتماعیِ تعیین‌کننده‌ی سلامت را یکپارچه می‌کند. ترخیص زودهنگام خطر بستری مجدد را افزایش می‌دهد؛ تأخیر در ترخیص ظرفیت \lr{ICU} را تحت فشار قرار می‌دهد و هزینه‌ها را بالا می‌برد.

در این پژوهش، قابلیت/امکان‌پذیری درمان به‌صورت عملیاتی به‌معنای پیش‌بینی یک پیامد بالینی قابل‌اندازه‌گیری بر اساس سیگنال‌های اولیه بیمار در \lr{ICU} تعریف می‌شود. این مفهوم به‌معنای توصیه مستقیم درمان، نسخه‌نویسی خودکار یا جایگزینی تصمیم پزشک نیست. بنابراین خروجی مدل، در بهترین حالت، یک نشانه‌ی کمّی برای پشتیبانی تصمیم و اولویت‌بندی بررسی بالینی است، نه دستور درمان.

مراقبت‌های بهداشتی مدرن به‌طور فزاینده‌ای به \emph{داده‌های تجمیع‌شده} (پایگاه‌های داده‌ی بیمارستانی، پرونده‌های \lr{ICU}، دستگاه‌های پوشیدنی و گزارش علائم از منابع متعدد) برای پشتیبانی از پیش‌بینی سلامت فردی تکیه دارد. امکان‌پذیری درمان، که در این پایان‌نامه به‌صورت ارزیابی دودوییِ دستیابی به ترخیص ایمن در بازه‌ی زمانی معقول و نبود بستری مجدد عملیاتی شده است، نقطه‌ی تصمیم کلیدی است. مدل‌های یادگیری ماشین که امکان‌پذیری را از داده‌های اولیه‌ی \lr{ICU} پیش‌بینی کنند می‌توانند به بالین‌گران در برنامه‌ریزی ترخیص کمک کنند، به‌شرط آنکه دقیق، قابل‌تفسیر، محدودیت‌مند و در شرایط واقعی اعتبارسنجی شده باشند.

%==============================================================================
\section{تحول دیجیتال در سلامت}
%==============================================================================

نظام سلامت در سال‌های اخیر دگرگونی دیجیتال گسترده‌ای را تجربه کرده است. گسترش پرونده‌ی الکترونیک سلامت (\lr{EHR})، ابزارهای پوشیدنی پایش سلامت، پزشکی از راه دور و استانداردهای تبادل داده مانند \lr{HL7 FHIR} باعث شده حجم و تنوع داده‌های بالینی به‌طور چشمگیری افزایش یابد. پیش‌بینی‌ها نشان می‌دهند بازار جهانی سلامت دیجیتال تا سال \pnum{2028} از \pnum{500} میلیارد دلار فراتر خواهد رفت؛ رشدی که نشان‌دهندۀ اهمیت روزافزون سامانه‌های پشتیبان تصمیم مبتنی بر داده است. این تحول، هم فرصت‌های تازه‌ای برای یادگیری ماشین ایجاد می‌کند و هم چالش‌های جدی دربارۀ کیفیت، حریم خصوصی و قابلیت تعمیم داده‌ها پیش می‌کشد.

\subsection{از سیلوهای داده تا هوش تجمیعی}

از نظر تاریخی، داده‌های بالینی در سیلوهای نهادی قرار داشت؛ هر بیمارستان سامانه‌ی \lr{EHR} خود را با طرحواره‌های اختصاصی، قراردادهای کدگذاری و کنترل‌های دسترسی نگه می‌داشت. ظهور تبادلات اطلاعات سلامت (\lr{HIE}) و شبکه‌های داده‌ی فدرال شروع به شکستن این موانع کرده است و پژوهش چندنهادی را بدون متمرکزسازی کامل داده ممکن می‌سازد. ابتکارهایی مانند \lr{PCORnet}، \lr{OHDSI} (\lr{Observational Health Data Sciences and Informatics}) و برنامه‌ی \lr{UK NHS Digital} پتانسیل هوش بالینی تجمیعی را نشان می‌دهند. بااین‌حال، این شبکه‌ها با چالش‌های پایدار روبه‌رو هستند: قالب‌های ناهمگون داده، کیفیت متغیر داده، امتناع نهادی از اشتراک پرونده‌های حساس و محدودیت‌های مقرراتی که در حوزه‌های قضایی مختلف متفاوت است.

این پایان‌نامه در همین بستر تعریف می‌شود. ایدۀ اصلی آن است که یادگیری فدرال با سازگاری کم‌نمونه ترکیب شود تا مدل بتواند از داده‌های بالینی توزیع‌شده بیاموزد، بدون آنکه داده‌های خام در یک مرکز تجمیع شوند. بازنمایی توکن‌های علائم نیز نقش یک لایۀ میانی استاندارد را دارد؛ یعنی داده‌های ناهمگون را به قالبی تبدیل می‌کند که برای مدل قابل استفاده باشد و در عین حال وابستگی مستقیم به قالب خام هر نهاد را کاهش دهد.

\subsection{وعده و خطر هوش مصنوعی بالینی}

هوش مصنوعی در مراقبت‌های بهداشتی در حوزه‌های مشخص نتایج چشمگیر به‌دست آورده است: رادیولوژی (مثلاً غربالگری رتینوپاتی دیابتی)، پاتولوژی (مثلاً تشخیص تومور) و پردازش زبان طبیعی (مثلاً خلاصه‌سازی یادداشت بالینی). بااین‌حال، انتقال پژوهش \lr{AI} به عمل بالینی کند مانده است. موانع کلیدی عبارت‌اند از: (\pnum{1}) \emph{شکاف‌های اعتبارسنجی}؛ مدل‌های آموزش‌دیده بر داده‌ی گذشته‌نگر ممکن است به استقرار آینده‌نگر تعمیم نیابند؛ (\pnum{2}) \emph{عدم قطعیت مقرراتی}؛ چارچوب در حال تکامل \lr{FDA} برای نرم‌افزار به‌عنوان دستگاه پزشکی (\lr{SaMD}) ابهام ایجاد می‌کند؛ (\pnum{3}) \emph{اعتماد بالین‌گر}؛ مدل‌های مبهم با مقاومت پزشکانی روبه‌رو می‌شوند که مسئولیت پیامدهای بیمار را بر عهده دارند؛ و (\pnum{4}) \emph{دسترسی به داده}؛ مقررات حریم خصوصی داده‌ی موجود برای توسعه‌ی مدل را محدود می‌کند.

این پایان‌نامه به‌طور مستقیم به دو مانع مهم می‌پردازد: محدودیت دسترسی به داده از طریق یادگیری فدرال، و شکاف اعتبارسنجی از طریق طراحی آزمایش و گزارش دقیق. نتیجۀ منفی اولیه، یعنی عملکرد ضعیف‌تر مدل کم‌نمونه نسبت به مدل‌های کلاسیک، پنهان نشده است؛ بلکه به‌عنوان نقطۀ شروع تحلیل مطرح می‌شود. سپس با چند دور بهبود مهندسی و مدل‌سازی، فاصله با سقف جدولی کاهش می‌یابد. این رویکرد با روح ارزیابی مسئولانۀ \lr{AI} در سلامت سازگار است، زیرا اعتماد بالینی زمانی شکل می‌گیرد که هم موفقیت‌ها و هم محدودیت‌ها شفاف گزارش شوند.

\subsection{امکان‌پذیری درمان در عصر مراقبت مبتنی بر ارزش}

گذار از مدل کارمزد به‌ازای خدمت به مدل‌های مراقبت مبتنی بر ارزش، برنامه‌ریزی ترخیص را به اولویت مالی و بالینی تبدیل کرده است. در برنامه‌های پرداخت تجمیعی (مثلاً \lr{CMS Bundled Payments for Care Improvement})، بیمارستان‌ها مسئولیت مالی بستری مجدد در \pnum{30} روز را بر عهده دارند. پیش‌بینی دقیق امکان‌پذیری درمان (اینکه آیا بیمار می‌تواند در بازه‌ی زمانی بالینی مناسب به ترخیص ایمن دست یابد) مستقیماً بر تخصیص منابع، مدیریت تخت و بازپرداخت اثر می‌گذارد. شناسایی زودهنگام بیمارانی که احتمال ترخیص ممکن را دارند، هماهنگی مراقبت پیشگیرانه را ممکن می‌سازد: ترتیب خدمات بهداشت خانگی، زمان‌بندی ارزیابی‌های توان‌بخشی و ارتباط با خانواده‌ها درباره‌ی جدول زمانی مورد انتظار.

بستر \lr{ICU} به‌ویژه سرنوشت‌ساز است؛ هزینه‌ی هر روز بستری در \lr{ICU} حدود \pnum{3} تا \pnum{5} برابر هزینه‌ی روزِ بستری در بخش عمومی است و ظرفیت \lr{ICU} در اوج موج‌ها محدودیت الزام‌آور است (همان‌طور که در همه‌گیری \lr{COVID-19} نشان داده شد). ابزار پشتیبان تصمیم که زمان‌بندی ترخیص را حتی یک روز بهبود دهد می‌تواند صرفه‌جویی هزینه‌ی قابل‌توجه و جریان بهتر بیمار را به همراه داشته باشد.

%==============================================================================
\section{انگیزه}
%==============================================================================

انگیزۀ اصلی این کار از یک تجربۀ عملی شکل گرفت. هنگام کار با پلتفرم \lr{Mayo Clinic} روشن شد که دسترسی به دادۀ سلامت، حتی برای هدف‌های پژوهشی، بسیار محدود است: در بسیاری موارد تنها شناسه‌ها و کدها در دسترس‌اند و پرونده‌های بالینی زیربنایی قابل استفاده نیستند. این تجربه پرسش اصلی پایان‌نامه را برجسته کرد: اگر داده را نتوان متمرکز یا آزادانه مبادله کرد، چگونه می‌توان سامانه‌های پیش‌بین قابل اعتماد ساخت؟ پاسخ این پایان‌نامه بررسی ترکیب یادگیری فدرال، سازگاری کم‌نمونه و بازنمایی‌های فشرده مانند \lr{Symptom Tokens} است.

بخش مراقبت ویژه (\lr{ICU}) یکی از غنی‌ترین محیط‌های داده‌ای و در عین حال بحرانی‌ترین محیط‌های تصمیم در مراقبت‌های مدرن است. بالین‌گران باید به‌طور مداوم پایداری بیمار، پیامدها و مقصد مناسب ترخیص را بر اساس مشاهدات اولیه‌ی محدود ارزیابی کنند. امکان‌پذیری درمان (پرسش دودوییِ اینکه آیا بیمار می‌تواند به‌صورت ایمن به منزل، توان‌بخشی یا سایر بسترهای غیرحاد ترخیص شود) به مقادیر آزمایشگاهی، علائم حیاتی و مسیرهای بالینی مشاهده‌شده در ساعات نخست بستری \lr{ICU} وابسته است.

علیرغم پیشرفت‌ها در \lr{EHR} و یادگیری ماشین، چند چالش بنیادین پابرجاست. نخست، \emph{کمیابی داده}: بسیاری از زیرگروه‌های بیماری نمونه‌های برچسب‌خورده‌ی اندک دارند و یادگیری نظارت‌شده‌ی استاندارد را دشوار می‌کند. دوم، \emph{حریم خصوصی}: داده‌ی بهداشتی به‌دلیل محدودیت‌های مقرراتی (مثلاً \lr{HIPAA}، \lr{GDPR}) را نمی‌توان بین نهادها متمرکز کرد. سوم، \emph{بازنمایی}: داده‌ی بالینی خام ناهمگون است (آزمایشگاه، علائم حیاتی، جمعیت‌شناسی) و برای مدل‌های پایین‌دست پیش‌پردازش دقیق می‌طلبد. چهارم، \emph{تعمیم}: مدل‌های آموزش‌دیده بر یک گروه بیماری ممکن است به گروه‌های دیگر منتقل نشوند.

یادگیری کم‌نمونه~\cite{snell2017prototypical,vinyals2016matching} این امکان را فراهم می‌کند که مدل با تعداد کمی نمونۀ پشتیبان برای هر دسته، بیماران جدید را طبقه‌بندی کند. از سوی دیگر، یادگیری فدرال~\cite{mcmahan2017communication} آموزش بر دادۀ توزیع‌شده را بدون انتقال نمونه‌های خام ممکن می‌سازد. ترکیب این دو رویکرد می‌تواند مدلی ایجاد کند که از دادۀ محدود بیاموزد، حریم خصوصی را حفظ کند و بین گروه‌های بیماری تعمیم یابد. این پایان‌نامه بررسی می‌کند که آیا چنین ترکیبی برای پیش‌بینی امکان‌پذیری درمان روی \lr{MIMIC\text{-}IV}~\cite{johnson2023mimic} عملی و قابل دفاع است یا خیر.

%==============================================================================
\section{بیان مسئله}
%==============================================================================

پیش‌بینی امکان‌پذیری درمان را به‌عنوان وظیفه‌ی طبقه‌بندی دودویی صورت‌بندی می‌کنیم. با مشاهدات بالینی بیمار از \pnum{24} ساعت نخست بستری در \lr{ICU}، پیش‌بینی می‌کنیم که آیا بیمار به ترخیص \emph{ممکن} دست خواهد یافت؛ تعریف‌شده به‌عنوان ترخیص به منزل، مراقبت خانگی یا توان‌بخشی ظرف \pnum{30} روز از پذیرش \lr{ICU}، بدون بستری مجدد ظرف \pnum{30} روز. برچسب از سوابق ترخیص و بستری مجدد \lr{MIMIC\text{-}IV} استخراج می‌شود.

به‌صورت رسمی، فضای بازنمایی بیماران را با $\mathcal{X}$ (مثلاً دنباله‌ی توکن‌های علائم یا بردارهای ویژگی جدولی) و برچسب دودویی را با $\mathcal{Y} = \{0, 1\}$ (\pnum{0} = غیرممکن، \pnum{1} = ممکن) نشان می‌دهیم. تابعی $f: \mathcal{X} \to \mathcal{Y}$ می‌جوییم که خطای پیش‌بینی را بر مجموعه‌ی آزمون نگه‌داشته‌شده کمینه کند، با قید اینکه $f$ ممکن است در بستر کم‌نمونه و فدرال آموزش داده شود.

چالش‌های کلیدی عبارت‌اند از:
\begin{enumerate}
\item \textbf{بستر کم‌نمونه}: در زمان استنتاج ممکن است برای هر دسته فقط $K$ نمونۀ برچسب‌خورده در اختیار باشد، برای مثال $K=5$. مدل باید بتواند از این نمونه‌های پشتیبان به بیماران پرس‌وجوی جدید تعمیم دهد.
\item \textbf{بستر فدرال}: دادۀ آموزشی میان $N$ گره، مانند گروه‌های بیماری، تقسیم شده است و دادۀ خام از گره‌های محلی خارج نمی‌شود. بنابراین آموزش باید از طریق تجمیع فدرال انجام شود.
\item \textbf{بازنمایی}: لازم است بازنمایی فشرده و معناداری از آزمایش‌ها و علائم حیاتی ساخته شود که هم برای مدل‌های کلاسیک و هم برای مدل‌های عصبی کم‌نمونه قابل استفاده باشد.
\end{enumerate}

\paragraph{چرا کم‌نمونه؟} در استقرار دنیای واقعی، زیرگروه‌های بیماری جدید یا نهادها ممکن است داده‌ی برچسب‌خورده‌ی محدود داشته باشند. یادگیری کم‌نمونه اجازه می‌دهد مدل تنها با $K$ نمونه برای هر دسته سازگار شود و نیاز به مجموعه‌های حجیم حاشیه‌نویسی‌شده در هر سایت را کاهش دهد. این امر به‌ویژه برای بیماری‌های نادر یا شرایط نوظهور اهمیت دارد.

\paragraph{چرا فدرال؟} داده‌ی بهداشتی به‌دلیل مقررات حریم خصوصی (\lr{HIPAA}، \lr{GDPR}) را نمی‌توان متمرکز کرد. یادگیری فدرال مدل‌ها را بر داده‌ی توزیع‌شده بدون جابه‌جایی پرونده‌های خام آموزش می‌دهد و همکاری بین‌نهادی را با حفظ حاکمیت داده ممکن می‌سازد.

\paragraph{چرا \lr{Symptom Tokens}؟} داده‌ی بالینی خام ناهمگون است: آزمایشگاه‌ها واحدهای مختلف دارند (\lr{mg/dL}، \lr{mmol/L})، علائم حیاتی دامنه‌های متفاوت دارند و اندازه‌گیری‌ها به‌صورت نامنظم نمونه‌برداری می‌شوند. بازنمایی مبتنی بر توکن‌های علائم این را به دنباله‌ای استاندارد تبدیل می‌کند که رمزگذارهای \lr{transformer} می‌توانند پردازش کنند. هر توکن هویت علامت، شدت (مقدار نرمال‌شده)، افست زمانی و سنجه (\lr{lab} در برابر \lr{vital}) را رمزگذاری می‌کند. این بازنمایی فشرده هم یادگیری اپیزودیک کم‌نمونه و هم مبانی جدولی کلاسیک را هنگام تخت‌کردن به بردار ویژگی پشتیبانی می‌کند.

%==============================================================================
\section{شکاف پژوهشی}
%==============================================================================

علیرغم پیشرفت چشمگیر در پیش‌بینی بالینی و یادگیری ماشین، شکاف آشکاری در تقاطع یادگیری کم‌نمونه، آموزش فدرال و داده‌ی بالینی جدولی وجود دارد. شکاف‌های مشخص زیر را که این پایان‌نامه را انگیزه می‌دهند، شناسایی می‌کنیم:

\subsection{شکاف \pnum{1}: یادگیری کم‌نمونه بر داده‌ی بالینی جدولی}

اکثریت قاطع پژوهش یادگیری کم‌نمونه بینایی ماشین (\lr{miniImageNet}، \lr{CUB-200}، \lr{Omniglot}) یا وظایف پردازش زبان طبیعی را هدف قرار می‌دهد. داده‌ی جدولی (سنجه‌ی غالب در پیش‌بینی بالینی مبتنی بر \lr{EHR}) در بستر کم‌نمونه به‌شدت کم‌بررسی مانده است. برخلاف تصاویر، ویژگی‌های جدولی ساختار فضایی که سوگیری القایی کانولوشنی بتواند از آن بهره برد ندارند. برخلاف متن، ویژگی‌های جدولی انواع و مقیاس‌های ناهمگونی دارند که در برابر توکن‌سازی استاندارد مقاومت می‌کنند. پرسش اینکه آیا متایادگیری مبتنی بر فاصله (شبکه‌های نمونه‌ای، \lr{prototypical networks}) می‌تواند بازنمایی مؤثر از ویژگی‌های بالینی ناهمگون با نمونه‌های محدود بیاموزد، همچنان باز است. در چارچوب مرور ادبیات این پژوهش، این پایان‌نامه این پرسش را بر دادۀ امکان‌پذیری درمان \lr{MIMIC\text{-}IV} به‌صورت تجربی و با خطوط پایۀ جدولی قوی بررسی می‌کند.

\subsection{شکاف \pnum{2}: یادگیری فدرال کم‌نمونه برای مراقبت‌های بهداشتی}

درحالی‌که یادگیری فدرال به مراقبت‌های بهداشتی اعمال شده (مثلاً \lr{FedHealth}، \lr{FL-QSAR} برای کشف دارو) و یادگیری کم‌نمونه به تصویربرداری پزشکی اعمال شده، ترکیب آموزش فدرال با متایادگیری کم‌نمونه بر داده‌ی بالینی ساختاریافته بررسی نشده است. این ترکیب از نظر عملی مهم است: سناریوهای استقرار واقعی اغلب هم کمیابی داده (نمونه‌های برچسب‌اندک در نهاد جدید) و هم قیود حریم خصوصی (عدم خروج داده از نهاد) را شامل می‌شوند. این پایان‌نامه بررسی می‌کند که آیا این دو پارادایم را می‌توان بدون افت فاجعه‌بار عملکرد به‌طور مؤثر ترکیب کرد.

\subsection{شکاف \pnum{3}: ارزیابی صادقانه‌ی روش‌های عصبی بر مجموعه‌داده‌های بالینی کوچک}

سوگیری انتشار در یادگیری ماشین معمولاً به سود نتایج مثبت عمل می‌کند. بسیاری از مقالات \lr{ML} بالینی فقط مقایسه‌های مطلوب را گزارش می‌کنند و شرایطی را که در آن روش‌های ساده‌تر برتری دارند کنار می‌گذارند. به همین دلیل، ارزیابی دقیق و صادقانه باید هم نتایج منفی اولیه و هم مسیر بهبود را نشان دهد~\cite{makel2012replications}. این پایان‌نامه با تعیین سقف جدولی قوی (\lr{SOTA} با \lr{AUROC} \textbf{\pnum{0.748}} در اجرای منجمد)، گزارش شکاف \lr{Exp1} (مدل پیشنهادی کم‌نمونه با \textbf{\pnum{0.611}} در برابر \lr{RF} با \textbf{\pnum{0.762}} روی همان تقسیم)، تحلیل دلیل عقب‌ماندن کم‌نمونه، و مستندسازی بهترین اجرای انباشته (\textbf{\pnum{0.746}} در برابر \textbf{\pnum{0.748}}) این شکاف گزارش‌دهی را تا حد زیادی پوشش می‌دهد. کمپین توسعۀ \lr{V1--V6} نیز در فصل~\ref{ch:results} به‌عنوان سابقۀ تاریخی تکامل روش آمده است.

\subsection{شکاف \pnum{4}: بازنمایی‌های بالینی مبتنی بر توکن‌های علائم}

کار موجود بر \lr{transformer}های بالینی (\lr{BEHRT}، \lr{Med-BERT}، \lr{ClinicalBERT}) بر یادداشت‌های بالینی (متن آزاد) یا کدهای تشخیصی عمل می‌کند. اعمال معماری‌های \lr{transformer} به اندازه‌گیری‌های بالینی کمی (آزمایشگاه، علائم حیاتی) به پل بازنمایی نیاز دارد؛ یعنی تبدیل اندازه‌گیری‌های پیوسته و ناهمگون به دنباله‌ی توکن. طرح توکن‌های علائم در این پژوهش این شکاف را پوشش می‌دهد و نگاشت اصولی از اندازه‌گیری‌های بالینی به ورودی‌های سازگار با \lr{transformer} فراهم می‌کند. اگرچه از نظر مفهومی به \lr{TabTransformer} و \lr{FT-Transformer} مرتبط است، رویکرد ما عناصر ویژه‌ی حوزه (هویت علامت، نرمال‌سازی بازه‌ی بالینی، رمزگذاری سنجه) را متناسب با بستر پیش‌بینی \lr{ICU} دربرمی‌گیرد.

%==============================================================================
\section{اصطلاحات و تعاریف}
%==============================================================================

در سراسر متن از اصطلاحات زیر استفاده می‌کنیم: \textbf{امکان‌پذیری درمان} پیامد دودویی دستیابی بیمار به ترخیص ایمن به منزل، توان‌بخشی یا مراقبت خانگی ظرف \pnum{30} روز بدون بستری مجدد را نشان می‌دهد. \textbf{\lr{Symptom Tokens}} بازنمایی‌های فشرده‌ی اندازه‌گیری آزمایشگاهی و علائم حیاتی هستند که هرکدام شناسه‌ی علامت، شدت، زمان و سنجه را رمزگذاری می‌کنند. \textbf{یادگیری کم‌نمونه} به متایادگیری از $K$ نمونه‌ی پشتیبان برای هر دسته اشاره دارد. \textbf{یادگیری فدرال} مدل‌ها را بر داده‌ی توزیع‌شده بدون متمرکزسازی نمونه‌های خام آموزش می‌دهد. \textbf{\lr{ETHOS}} رمزگذار \lr{Efficient Transformer for Health Outcome prediction} است~\cite{sharafi2023ethos}.

%==============================================================================
\section{مرور کلی روش پژوهش}
%==============================================================================

پژوهش ما از رویکرد علم طراحی پیروی می‌کند: یک چارچوب ترکیبی شامل \lr{Symptom Tokens}، \lr{ETHOS}، یادگیری کم‌نمونه و آموزش فدرال طراحی، پیاده‌سازی و به‌صورت تجربی ارزیابی می‌شود. برای تکرارپذیری، از کوهورت \pnum{2,000} بیمار \lr{MIMIC\text{-}IV} با تقسیم‌های ثابت استفاده می‌کنیم. ارزیابی شامل پنج بخش است: (\pnum{1}) مقایسه با مدل‌های کلاسیک مانند \lr{RF} و \lr{XGBoost}؛ (\pnum{2}) حذف تدریجی اجزا؛ (\pnum{3}) تحلیل حساسیت برای \lr{K-shot} و ابرپارامترها؛ (\pnum{4}) مقایسۀ فدرال و متمرکز؛ و (\pnum{5}) تعمیم بین بیماری‌ها. هم نتایج مثبت و هم نتایج منفی گزارش می‌شوند تا تصویر علمی کامل‌تری ارائه شود.

%==============================================================================
\section{اهداف پژوهش و مفروضات}
%==============================================================================

هدف اصلی این پایان‌نامه طراحی و ارزیابی چارچوبی پژوهشی برای پیش‌بینی امکان‌پذیری درمان از داده‌ی بالینی تجمیع‌شده و محدود است. چارچوب باید الگوهای مرتبط با پیامد ترخیص/درمان را بدون اتکا به جابه‌جایی مستقیم داده‌ی خام بیمار بررسی کند. اهداف مشخص عبارت‌اند از: (\pnum{1}) ارائه و ارزیابی مدلی مبتنی بر یادگیری کم‌نمونه برای پیش‌بینی امکان‌پذیری درمان؛ (\pnum{2}) استفاده از \lr{Symptom Tokens} به‌عنوان واحدهای استاندارد داده برای یادگیری مؤثر؛ (\pnum{3}) ترکیب معماری \lr{ETHOS} با اصول یادگیری کم‌نمونه و فدرال برای سنجش سودمندی آن در محیط‌های چندمنبعه؛ و (\pnum{4}) فراهم‌کردن بستری پژوهشی برای تحلیل سلامت فردی با کاهش نیاز به تمرکز داده‌ی خام.

بر اساس ادبیات و چارچوب نظری، مفروضات آزمون‌پذیر زیر مطرح می‌شوند: (\pnum{1}) مدل‌های کم‌نمونه ممکن است در شرایط داده‌ی محدود به پیش‌بینی امکان‌پذیری درمان کمک کنند، اما باید با خطوط پایۀ جدولی قوی مقایسه شوند؛ (\pnum{2}) یادگیری فدرال می‌تواند نیاز به انتقال داده‌ی خام را کاهش دهد، هرچند به‌تنهایی تضمین کامل حریم خصوصی نیست؛ (\pnum{3}) ترکیب \lr{Symptom Tokens} با معماری \lr{ETHOS} ممکن است بازنمایی مفیدی فراهم کند؛ (\pnum{4}) الگوهای پنهان پیامد درمان/ترخیص باید با تحلیل خطا و اهمیت ویژگی تفسیر شوند؛ و (\pnum{5}) داده‌ی چندسنجه‌ای در مقایسه با داده‌ی تک‌سنجه‌ای می‌تواند اطلاعات بیشتری برای پیش‌بینی فراهم کند. نتایج فصل~\ref{ch:results} نشان می‌دهد برخی از این مفروضات تنها در حالت ترکیبی و انباشتی پشتیبانی می‌شوند، نه در مدل کم‌نمونه‌ی خام.

\subsection{تطبیق با پروپوزال اولیه}

پروپوزال اولیه چارچوبی با چهار رکن اصلی پیشنهاد می‌کرد: یادگیری کم‌نمونه، یادگیری فدرال، رمزگذار \lr{ETHOS} و بازنمایی \lr{Symptom Tokens} برای پیش‌بینی قابلیت درمان از داده‌های بالینی تجمیع‌شده و محدود. پایان‌نامه حاضر همان مسیر اصلی را حفظ می‌کند، اما برای اینکه ارزیابی علمی قابل بازتولید و قابل دفاع باشد، دامنه‌ی اجرا را به کوهورت ناشناس‌شده و عمومی \lr{MIMIC\text{-}IV} محدود می‌کند. به‌دلیل نبود دسترسی به داده‌ی واقعی چندمرکزی، جریان پوشیدنی و داده‌ی زنده در چارچوب زمانی این پژوهش کارشناسی ارشد، «داده‌ی چندمنبع» در این نسخه به‌صورت آزمایشگاه‌ها، علائم حیاتی، جمعیت‌شناسی و ویژگی‌های مهندسی‌شده از \lr{ICU} عملیاتی شده است. بنابراین داده‌ی ابزارهای پوشیدنی، واسط کاربری تولیدی و استقرار چندسایتی واقعی در این پایان‌نامه پیاده‌سازی نشده‌اند و به‌عنوان مسیرهای آینده و الزامات استقرار واقعی مطرح می‌شوند.

\begin{table}[htbp]
\centering
\caption{تطبیق اجزای پروپوزال با اجرای نهایی پایان‌نامه}
\footnotesize
\setlength{\tabcolsep}{1.8pt}
\renewcommand{\arraystretch}{1.12}
\begin{tabularx}{\linewidth}{@{}R{2.25cm}R{2.85cm}C{1.35cm}R{3.05cm}Y@{}}
\toprule
\textbf{جزء در پروپوزال} & \textbf{اجرای نهایی در پایان‌نامه} & \textbf{سطح پوشش} & \textbf{توضیح دامنه} & \textbf{نکتۀ دفاعی} \\
\midrule
یادگیری کم‌نمونه & شبکه‌های نمونه‌اولیه، تحلیل $K$-شات، \lr{Proto-MAML} و مسیرهای \lr{V1--V6} & کامل & نتایج مثبت و محدودیت‌ها هر دو گزارش شده‌اند & ادعا نمی‌شود مدل خام از \lr{RF} بهتر است \\
یادگیری فدرال & شبیه‌سازی \lr{FedAvg} روی گره‌های بیماری & جزئی & داده‌ی خام در شبیه‌سازی جابه‌جا نمی‌شود؛ استقرار واقعی انجام نشده است & هزینۀ فدرال در این تنظیم حدود \pnum{0.005} \lr{AUROC} است \\
\lr{ETHOS} و \lr{Symptom Tokens} & توکن‌سازی \pnum{15} آزمایشگاه، \pnum{7} علامت حیاتی و ویژگی‌های مهندسی‌شده، سپس رمزگذاری با \lr{ETHOS} & کامل & بازنمایی پیشنهادی روی \lr{MIMIC\text{-}IV} بازتولیدپذیر شده است & نقش آن بیشتر در مسیر ترکیبی/انباشتی دفاع‌پذیر است \\
داده‌ی تجمیع‌شده/چندسنجه‌ای & آزمایشگاه، علائم حیاتی، جمعیت‌شناسی، گره بیماری و ویژگی‌های بالینی مشتق‌شده & جزئی & ابزار پوشیدنی و داده‌ی زنده در دسترس نبود؛ ادعای چندمنبع به منابع درون \lr{EHR} محدود است & این یک تصمیم دامنه‌بندی برای پژوهش کارشناسی ارشد است \\
پیش‌بینی قابلیت درمان & پیامد جانشین ترخیص/بستری مجدد از \lr{MIMIC\text{-}IV} & کامل با محدودیت & برچسب، جانشین عملیاتی است نه حقیقت کامل بالینی & مدل پشتیبان تصمیم است، نه نسخه‌نویس خودکار \\
مقایسه با روش‌های سنتی & \lr{LR}، \lr{RF}، \lr{XGBoost}، \lr{LightGBM}، \lr{CatBoost} و مدل‌های عصبی & کامل & سقف جدولی برای تفسیر منصفانه گزارش شده است & مدل‌های کلاسیک روی داده‌ی جدولی بسیار قوی‌اند \\
واسط کاربری و نمایش نتایج & جداول، نمودارها، \lr{SHAP} و گزارش‌های قابل بازتولید & آینده & واسط کاربری تولیدی ساخته نشده است & تمرکز پایان‌نامه اعتبار علمی مدل و نتایج تجربی است \\
ابزارهای پیشنهادی مانند \lr{PySyft}/\lr{Flutter} & پیاده‌سازی پژوهشی با \lr{Python}، \lr{PyTorch}، \lr{scikit-learn} و شبیه‌سازی سفارشی \lr{FedAvg} & جایگزین موجه & انتخاب ابزارها بر اساس بازتولیدپذیری و دسترسی واقعی داده انجام شده است & ابزار عوض شده، پرسش پژوهش حفظ شده است \\
\bottomrule
\end{tabularx}
\label{tab:proposal_alignment}
\end{table}

%==============================================================================
\section{پرسش‌های پژوهش}
%==============================================================================

این پایان‌نامه به پرسش‌های پژوهشی زیر می‌پردازد:

\begin{enumerate}
\item \textbf{\lr{RQ1}:} آیا متایادگیری کم‌نمونه می‌تواند در پیش‌بینی امکان‌پذیری درمان از \lr{MIMIC\text{-}IV} با مبانی جدولی کلاسیک برابری کند یا از آن پیشی بگیرد؟

\item \textbf{\lr{RQ2}:} تعداد شات‌های پشتیبان $K$ چگونه بر عملکرد کم‌نمونه اثر می‌گذارد؟ آیا $K$ بهینه‌ای وجود دارد؟

\item \textbf{\lr{RQ3}:} هزینه‌ی عملکردی یادگیری فدرال در مقایسه با آموزش متمرکز چیست؟ آیا این هزینه ناچیز است؟

\item \textbf{\lr{RQ4}:} کدام اجزای معماری (توکن‌های علائم، وزن‌دهی شدت، افت زمانی و غیره) بیشترین سهم را در مدل پیشنهادی دارند؟

\item \textbf{\lr{RQ5}:} آیا مدل بین گروه‌های بیماری تعمیم می‌یابد؟ بر گره‌های دیده‌نشده (مثلاً کلیوی، دیابت) هنگام آموزش بر سایرین (سپسیس، قلبی، تنفسی) چگونه عمل می‌کند؟
\end{enumerate}

به این پرسش‌ها به‌صورت تجربی با نمونه‌ی \pnum{2,000} بیمار از \lr{MIMIC\text{-}IV v3.1}، با تقسیم‌های ثابت آموزش/اعتبارسنجی/آزمون و پروتکل‌های ارزیابی دقیق پاسخ می‌دهیم.

%==============================================================================
\section{نوآوری‌ها و مشارکت‌ها}
%==============================================================================

این پایان‌نامه مشارکت‌های زیر را دارد:

\begin{enumerate}
\item \textbf{بررسی یکپارچۀ یادگیری کم‌نمونه، توکن‌سازی علائم و فدرال برای امکان‌پذیری درمان.} در چارچوب مرور ادبیات این پژوهش، ترکیب توکن‌سازی علائم، متایادگیری نمونه‌ای و آموزش فدرال برای پیش‌بینی امکان‌پذیری درمان بر \lr{MIMIC\text{-}IV} کمتر به‌صورت یکپارچه بررسی شده است. این مشارکت شکاف بین جامعه‌ی یادگیری کم‌نمونه (عمدتاً متمرکز بر بینایی و \lr{NLP}) و جامعه‌ی انفورماتیک بالینی (عمدتاً متمرکز بر یادگیری نظارت‌شده‌ی متمرکز) را به‌صورت پژوهشی بررسی می‌کند. معماری به‌گونه‌ای ماژولار طراحی شده که رمزگذار، سر کم‌نمونه و پروتکل فدرال هرکدام به‌طور مستقل قابل جایگزینی‌اند و حذف اجزا و گسترش آینده را تسهیل می‌کنند.

\item \textbf{\lr{Symptom Tokens} به‌عنوان بازنمایی بالینی.} بازنمایی فشرده‌ی مبتنی بر توکن معرفی می‌کنیم: \pnum{15} مورد آزمایشگاهی و \pnum{7} علامت حیاتی، نرمال‌شده به $[0,1]$، تقویت‌شده با \lr{shock index}، پراکسی \lr{SOFA} و شمار آزمایشگاه‌های غیرطبیعی. هر توکن آرایه‌ی \lr{[symptom\_id, intensity, time\_delta, modality]} را رمزگذاری می‌کند. این بازنمایی هم رمزگذارهای سبک \lr{transformer} و هم مدل‌های جدولی کلاسیک را پشتیبانی می‌کند. برخلاف رویکردهای پیشین توکن‌سازی بالینی که بر یادداشت‌های متن آزاد یا کدهای تشخیصی عمل می‌کنند، طرح ما \emph{اندازه‌گیری‌های کمی} را به واژگان گسسته‌ی توکن با مقادیر شدت پیوسته تبدیل می‌کند و هم هویت دسته‌ای و هم بزرگی عددی هر مشاهده‌ی بالینی را حفظ می‌کند.

\item \textbf{سقف تجربی و بستن تدریجی شکاف.} سقف جدولی قوی، یعنی \textbf{\lr{AUROC} \pnum{0.748}} در اجرای منجمد \lr{SOTA} و \pnum{0.760} برای \lr{RF} با آستانۀ تنظیم‌شده در جدول کمپین \lr{V}، همراه با نقطۀ آغاز تاریخی کم‌نمونه (\textbf{\pnum{0.619}}) گزارش می‌شود. اجرای منجمد انباشتۀ کم‌نمونه + کلاسیک در مخزن به \textbf{\pnum{0.746}} \cnum{[0.699, 0.797]} می‌رسد؛ یعنی فقط حدود \textbf{\pnum{0.002}} پایین‌تر از برآورد نقطه‌ای \lr{SOTA}. بهترین بذر \lr{V5} (\textbf{\pnum{0.737}}) و انباشت \lr{V4} (\textbf{\pnum{0.732}}) نیز به‌عنوان نقاط عطف توسعه مستند شده‌اند.

\item \textbf{معماری آماده‌ی فدرال با هزینه‌ی شبیه‌سازی‌شده‌ی کوچک.} در آزمایش گزارش‌شده‌ی فدرال در برابر متمرکز، شکاف \lr{AUROC} حدود \textbf{\pnum{0.005}} است؛ آموزش حافظ حریم خصوصی در این شبیه‌سازی با هزینه‌ی کوچک عملی است. استقرار چندسایتی واقعی همچنان نیازمند اعتبارسنجی عملیاتی است.

\item \textbf{حذف سیستماتیک اجزا و تحلیل حساسیت جامع.} به‌صورت سیستماتیک اجزا (توکن‌های علائم، وزن‌دهی شدت، افت زمانی، سر کم‌نمونه) را حذف می‌کنیم و $K \in \{1,3,5,7,10,15,20\}$ را تغییر می‌دهیم. جداول و تفاسیر مفصل برای کار آینده ارائه می‌دهیم. حذف اجزا نشان می‌دهد وزن‌دهی شدت بحرانی‌ترین جز است ($\Delta$\lr{AUROC} $= -0.075$ هنگام حذف)، درحالی‌که سر کم‌نمونه را می‌توان با سر نظارت‌شده با افت کمینه جایگزین کرد ($\Delta$\lr{AUROC} $= -0.009$). این الگو نشان می‌دهد پروتکل آموزش اپیزودیک محرک اصلی عملکرد نیست.

\item \textbf{متایادگیر انباشته و آموزش پیشرفته.} انباشتگر کمپین \lr{V4} به \lr{AUROC} \pnum{0.732} می‌رسد و پیکربندی منجمد \texttt{\lr{exp\_fewshot\_best}} مقدار \textbf{\pnum{0.746}} را ثبت می‌کند. ترکیب سازگاری به‌سبک \lr{Proto-MAML}، تقطیر دانش، پیش‌آموزش تضادی نظارت‌شده و مهندسی ویژگی (\lr{22$\to$132} ویژگی) نشان می‌دهد روش‌های کم‌نمونه، در صورت طراحی سیستماتیک، می‌توانند به \lr{SOTA} جدولی این کوهورت نزدیک شوند.

\item \textbf{نوآوری‌های معماری برای یادگیری کم‌نمونه جدولی.} \lr{V5} توجه ویژگی دروازه‌دار (\lr{Gated Feature Attention}) (انتخاب ویژگی تطبیقی به‌ازای نمونه الهام‌گرفته از \lr{TabNet}) و شبکه‌های نمونه‌ای توجه متقاطع (\lr{Cross-Attention Prototypical Networks}) (شباهت وابسته به زمینه) را معرفی می‌کند. \lr{V6} مجموعه‌های عمیق ناهمگون ترکیب‌کننده‌ی \lr{FT-Transformer}، \lr{GFA-Transformer} و \lr{MLP-Mixer} را با پیش‌بینی مطابق (\lr{conformal prediction}) برای سنجش عدم‌قطعیت ارزیابی می‌کند. این آزمایش‌ها نشان می‌دهد مکانیسم‌های توجه برای یادگیری کم‌نمونه جدولی ضروری‌اند و انباشت کالیبره‌شده‌ی ساده بر متایادگیرهای پیچیده در مقیاس داده‌ی کوچک برتری دارد.
\end{enumerate}

%==============================================================================
\section{ساختار پایان‌نامه}
%==============================================================================

ادامه‌ی این پایان‌نامه به‌صورت زیر سازمان‌دهی شده است.

\textbf{فصل~\ref{ch:background}} کارهای مرتبط را مرور می‌کند: یادگیری کم‌نمونه، یادگیری فدرال، بازنمایی داده‌ی بالینی و مجموعه‌داده‌ی \lr{MIMIC\text{-}IV}. درباره‌ی شبکه‌های نمونه‌ای، شبکه‌های تطبیق، \lr{MAML}، \lr{FedAvg} و کاربردهای پیشین در مراقبت‌های بهداشتی بحث می‌کنیم.

\textbf{فصل~\ref{ch:methods}} روش‌شناسی را به‌تفصیل شرح می‌دهد: استخراج مجموعه‌داده، توکن‌سازی علائم، معماری رمزگذار \lr{ETHOS}، پروتکل کم‌نمونه، آموزش فدرال و معیارهای ارزیابی. صورت‌بندی‌های ریاضی برای همه‌ی اجزا ارائه می‌شود.

\textbf{فصل~\ref{ch:experiments}} بستر آزمایشی را ارائه می‌کند: تقسیم داده، ابرپارامترها، مدل‌های مبنایی (رگرسیون لجستیک، جنگل تصادفی، \lr{XGBoost}، \lr{LightGBM}، \lr{CatBoost}) و پیکربندی \lr{ETHOS} همراه با یادگیری کم‌نمونه. آزمایش‌های \prange{1}{5} (مقایسه‌ی اصلی، حساسیت \lr{K}-شات، فدرال در برابر متمرکز، حذف اجزا، تعمیم بین بیماری‌ها) را شرح می‌دهیم.

\textbf{فصل~\ref{ch:results}} نتایج را گزارش می‌کند: \lr{SOTA} جدولی منجمد (\textbf{\pnum{0.748}})، کم‌نمونه انباشته‌ی منجمد (\textbf{\pnum{0.746}})، \lr{Exp1} و آزمایش‌های \prange{2}{5}، جدول توسعه‌ی \lr{V1--V6} (نقاط عطف تاریخی: \pnum{0.732}، \pnum{0.737}، \pnum{0.760})، راهنمای آماری، قابل‌تفسیر بودن و تجسم‌های اضافی.

\textbf{فصل~\ref{ch:data_analysis}} تحلیل جامع داده را ارائه می‌کند: طرحواره‌ی \lr{MIMIC\text{-}IV}، جمعیت‌شناسی کوهورت، توزیع گره‌ی بیماری، همبستگی‌های آزمایشگاهی، پراکسی \lr{SOFA}، عملکرد طبقه‌بندی‌شده بر اساس سن، تحلیل خطا، نهفته‌سازی \lr{t-SNE}، تجسم نمونه‌ها و خط زمانی بیماران.

\textbf{فصل~\ref{ch:discussion}} نتایج را تفسیر می‌کند، محدودیت‌های متایادگیری کم‌نمونه را بررسی می‌کند و پیامدهای عملی و مسیرهای آیندۀ پژوهش را مرور می‌کند.

\textbf{فصل~\ref{ch:conclusion}} یافته‌ها و مشارکت‌های اصلی را خلاصه می‌کند.

\textbf{پیوست‌ها} جداول ابرپارامتر، پرسش‌های \lr{SQL} \lr{MIMIC\text{-}IV}، قطعه‌کدها، جداول کامل نتایج و استنتاج‌های ریاضی را فراهم می‌کنند.

%==============================================================================
\section{بستر مقرراتی و اخلاقی}
%==============================================================================

داده‌ی بهداشتی تابع مقررات سخت‌گیرانه است. در ایالات متحده، قانون قابلیت انتقال و پاسخگویی بیمه‌ی سلامت (\lr{HIPAA}) استفاده و افشای اطلاعات سلامت محافظت‌شده (\lr{PHI}) را نظارت می‌کند. در اتحادیه‌ی اروپا، مقررات عمومی حفاظت از داده (\lr{GDPR}) قیود مشابهی اعمال می‌کند. این مقررات توانایی متمرکزسازی داده بین نهادها را محدود می‌کند و یادگیری فدرال و رویکردهای حافظ حریم خصوصی را انگیزه می‌بخشد. کار ما از داده‌ی ناشناس‌شده‌ی \lr{MIMIC\text{-}IV} استفاده می‌کند که تحت رویه‌های سخت‌گیرانه‌ی ناشناس‌سازی قرار گرفته است. برای استقرار، هر مدلی باید با مقررات محلی و تأیید کمیته‌ی اخلاق پژوهش (\lr{IRB}) سازگار باشد.

%==============================================================================
\section{تحلیل ذی‌نفعان}
%==============================================================================

چند ذی‌نفع از پیش‌بینی امکان‌پذیری درمان بهره می‌برند:

\begin{itemize}
\item \textbf{بالین‌گران}: پشتیبان تصمیم برای برنامه‌ریزی ترخیص؛ شناسایی زودهنگام بیمارانی که ممکن است بستری طولانی یا انتقال نیاز داشته باشند.
\item \textbf{مدیران بیمارستان}: تخصیص منابع، مدیریت تخت و برنامه‌ریزی ظرفیت.
\item \textbf{بیماران و خانواده‌ها}: شفافیت درباره‌ی مسیر ترخیص مورد انتظار؛ کاهش عدم‌قطعیت.
\item \textbf{پرداخت‌کنندگان}: مهار هزینه از طریق زمان‌بندی مناسب ترخیص؛ کاهش بستری مجدد.
\item \textbf{پژوهشگران}: روش‌شناسی برای پیش‌بینی بالینی کم‌نمونه و فدرال؛ معیارهایی برای کار آینده.
\end{itemize}

رویکرد ما قابل‌تفسیر بودن (\lr{SHAP}، اهمیت ویژگی) و حریم خصوصی (یادگیری فدرال) را اولویت می‌دهد تا نگرانی‌های بالین‌گر و مقرراتی را پوشش دهد.

%==============================================================================
\section{خلاصه‌ی فصل}
%==============================================================================

این فصل پایان‌نامه را معرفی کرد: بستر بالینی تصمیم‌گیری در \lr{ICU} و برنامه‌ریزی ترخیص، انگیزه ناشی از محدودیت‌های دسترسی به داده در دنیای واقعی (تجربه‌ی \lr{Mayo Clinic})، بیان رسمی مسئله، پرسش‌های پژوهش و مشارکت‌ها. امکان‌پذیری درمان را به‌عنوان وظیفه‌ی طبقه‌بندی دودویی بر داده‌ی \pnum{24} ساعت نخست \lr{ICU} تعریف کردیم، چالش‌های کلیدی (کم‌نمونه، فدرال، بازنمایی) را شناسایی کردیم و ساختار پایان‌نامه را ترسیم کردیم. فصل بعد مبانی نظری را مرور می‌کند.

%==============================================================================
\section{دامنه و محدودیت‌ها}
%==============================================================================

این پایان‌نامه بر نمونه‌ی \pnum{2,000} بیمار از \lr{MIMIC\text{-}IV v3.1} متمرکز است. نتایج ممکن است به کوهورت‌های بزرگ‌تر یا نهادهای دیگر تعمیم نیابد. از برچسب دودویی امکان‌پذیری مشتق‌شده از محل ترخیص و بستری مجدد استفاده می‌کنیم؛ تعاریف جایگزین (مثلاً وضعیت عملکردی، زمان تا ترخیص) می‌تواند به نتیجه‌گیری‌های متفاوت منجر شود. استنتاج علّی یا قابل‌تفسیر بودن را عمیق بررسی نمی‌کنیم. بستر فدرال بخش‌بندی مبتنی بر بیماری را شبیه‌سازی می‌کند؛ یادگیری فدرال دنیای واقعی ممکن است با چالش‌های اضافی (هزینه‌ی ارتباط، خرابی گره، کلاینت‌های بیزانسی) روبه‌رو شود.

%==============================================================================
\section{اهمیت در بستر پزشکی دقیق}
%==============================================================================

این کار به دستور کار گسترده‌تر پزشکی دقیق (تنظیم درمان پزشکی بر ویژگی‌های فردی هر بیمار) می‌افزاید. اگرچه پزشکی دقیق سنتی بر پروفایل‌سازی ژنومی و درمان‌های هدفمند تمرکز دارد، مفهوم به‌طور طبیعی به \emph{آینده‌نگری دقیق} گسترش می‌یابد: پیش‌بینی مسیر فردی بیمار بر اساس پروفایل بالینی ویژه به‌جای میانگین‌های جمعیتی.

\subsection{از میانگین‌های جمعیتی تا پیش‌بینی فردی}

برنامه‌ریزی ترخیص سنتی بر قواعد تجربی بالینی و هنجارهای نهادی تکیه دارد: «بیشتر بیماران سپسیس با بهبود لاکتات تا روز پنجم ترخیص می‌شوند.» چنین قواعدی میانگین‌های جمعیتی‌اند که تغییرپذیری فردی را نادیده می‌گیرند. مسیر ترخیص بیمار \pnum{45} ساله مبتلا به سپسیس با عملکرد کلیوی طبیعی و حمایت اجتماعی قوی بسیار متفاوت از بیمار \pnum{78} ساله با سپسیس، بیماری مزمن کلیوی و مراقبت خانگی محدود است. مدل ما تلاش می‌کند این تفاوت‌های فردی را با یادگیری از ترکیب مشخص مقادیر آزمایشگاهی، علائم حیاتی و ویژگی‌های بالینی مشاهده‌شده در \pnum{24} ساعت نخست بستری \lr{ICU} هر بیمار بازتاب دهد.

\subsection{سازگاری کم‌نمونه به‌عنوان شخصی‌سازی}

چارچوب کم‌نمونه چارچوب طبیعی برای شخصی‌سازی ارائه می‌دهد: با تعداد اندکی بیمار مرجع (مجموعه‌ی پشتیبان) که بالینی شبیه بیمار پرسشی‌اند، مدل پیش‌بینی‌های خود را سازگار می‌کند. این امر مشابه استدلال بالینی است؛ پزشکی که بیمار جدید را ارزیابی می‌کند موارد مشابه را به یاد می‌آورد و از آن تجربیات برای آینده‌نگری بهره می‌گیرد. شبکه‌های نمونه‌ای این را با محاسبه‌ی نمونه‌های دسته (مرکزهای بیماران مشابه) و طبقه‌بندی بیماران پرسشی بر اساس نزدیکی رسمی می‌کنند. اگرچه نتایج ما نشان می‌دهد این رسمی‌سازی هنوز به توان پیش‌بینی مدل‌های کلاسیک نمی‌رسد، هم‌ترازی مفهومی با استدلال بالینی نشان می‌دهد که با بازنمایی‌های بهتر و مجموعه‌های مرجع بزرگ‌تر، شخصی‌سازی کم‌نمونه می‌تواند ابزار بالینی قدرتمندی شود.

\subsection{تعمیم بین‌نهادی}

پزشکی دقیق به مدل‌هایی نیاز دارد که بین نهادها تعمیم یابند؛ هر نهاد جمعیت‌شناسی بیمار، شیوه‌های بالینی و استانداردهای مستندسازی خود را دارد. یادگیری فدرال آموزش بین‌نهادی بدون اشتراک داده را ممکن می‌سازد، اما چالش فراتر از حریم خصوصی است: مدل‌ها باید بازنمایی‌هایی بیاموزند که نسبت به ویژگی‌های نهادی ثابت بمانند و در عین حال به تغییرات بالینی معنادار حساس بمانند. بخش‌بندی گره-بیماری ناهمگونی نهادی را با ایجاد توزیع‌های داده‌ی غیر \lr{IID} شبیه‌سازی می‌کند. آزمایش‌های تعمیم بین بیماری (فصل~\ref{ch:results}) بررسی می‌کنند که آیا مدل می‌تواند به جمعیت‌های بیمار دیده‌نشده سازگار شود؛ نیاز کلیدی برای استقرار در بسترهای متنوع مراقبت‌های بهداشتی.
